\documentclass[12pt,a4paper]{elegantbook}
\title{CookBook}
\subtitle{}
\author{Cookbook Generator}
\institute{}
\date{\today}
\version{1.0}

% 数学 / 字体 / 中文
\usepackage{amsmath,amssymb}
\usepackage{fontspec}

% 表格
\usepackage{booktabs}
\usepackage{tabularx}
\usepackage{xltabular}
\usepackage{longtable}
\usepackage{array}
\usepackage{makecell}
\usepackage{multirow}
\usepackage{threeparttable}
\usepackage{pdflscape}
\usepackage{needspace}
\renewcommand{\arraystretch}{1.25}
\setlength{\tabcolsep}{5pt}

% 代码块
\usepackage{listings}
\usepackage[most]{tcolorbox}
\lstdefinestyle{cookbookcode}{
  basicstyle=\ttfamily\small,
  breaklines=true,
  breakatwhitespace=true,
  columns=fullflexible,
  keepspaces=true,
  showstringspaces=false,
  frame=none,
  tabsize=2,
  extendedchars=true,
  literate={-}{{-}}1
}
\newtcblisting{codeblock}[2][]{
  listing only,
  listing style=cookbookcode,
  title={#2},
  breakable,
  enhanced,
  colback=gray!3,
  colframe=gray!45,
  boxrule=0.4pt,
  arc=2mm,
  left=1.5mm, right=1.5mm, top=1mm, bottom=1mm,
  fonttitle=\bfseries\small,
  #1
}

% 超链接
\usepackage{hyperref}
\hypersetup{colorlinks=true, linkcolor=blue, urlcolor=blue}

\begin{document}
\maketitle
\tableofcontents

% Auto-generated from appendix-a.mdx — edit appendix-a.mdx then re-run mdx2tex.mjs
\chapter{附录 A · functools 常用装饰器速查}
\label{ch:appendix-a}

\section*{附录 A · functools 常用装饰器速查}
Python 标准库 \textbackslash\{\}texttt\{functools\} 提供了几个常用装饰器/工具，自己写装饰器时几乎都会用到。

\section{速查表}
\begin{table}[htbp]
\centering
\begin{tabularx}{\textwidth}{X X X}
\toprule
装饰器 / 函数 & 用途 & 典型写法 \\
\midrule
\textbackslash\{\}texttt\{@functools.wraps(func)\} & 拷贝原函数 \textbackslash\{\}texttt\{\_\_name\_\_\}/\textbackslash\{\}texttt\{\_\_doc\_\_\} & wrapper 定义处加上 \\
\textbackslash\{\}texttt\{@functools.lru\_cache(maxsize=N)\} & 记忆化（LRU 缓存） & 递归、纯函数 \\
\textbackslash\{\}texttt\{@functools.cache\} & Python 3.9+，等价无上限 LRU & 递归、纯函数 \\
\textbackslash\{\}texttt\{@functools.singledispatch\} & 按第一个参数类型分派（泛函数） & 多态工具函数 \\
\textbackslash\{\}texttt\{@functools.total\_ordering\} & 补全比较运算符 & 只写 \textbackslash\{\}texttt\{\_\_eq\_\_\}+\textbackslash\{\}texttt\{\_\_lt\_\_\} 即可 \\
\textbackslash\{\}texttt\{functools.partial(fn, *args)\} & 偏函数：冻结部分参数 & 回调、工厂 \\
\textbackslash\{\}texttt\{functools.reduce(fn, iter)\} & 累积归约 & 替代简单 for-累加 \\
\bottomrule
\end{tabularx}
\end{table}

\section{使用示例}

\begin{codeblock}{python}
# file: examples/functools_demo.py
import functools

@functools.lru_cache(maxsize=128)
def fib(n: int) -> int:
    """记忆化斐波那契，复杂度从 O(2^n) 降到 O(n)。"""
    if n < 2:
        return n
    return fib(n - 1) + fib(n - 2)

print(fib(40))  # 瞬间出结果；没有 lru_cache 要算很久

@functools.total_ordering
class Version:
    def __init__(self, major: int, minor: int):
        self.major, self.minor = major, minor
    def __eq__(self, other): return (self.major, self.minor) == (other.major, other.minor)
    def __lt__(self, other): return (self.major, self.minor) < (other.major, other.minor)
    # 有了 total_ordering，>=、<=、> 自动生成

print(Version(1, 2) <= Version(1, 10))  # True
\end{codeblock}

\section{延伸阅读}
\begin{itemize}
\item Python 官方文档：https://docs.python.org/3/library/functools.html
\item Recipe 2-1 用到 \textbackslash\{\}texttt\{functools.wraps\}：见 第 2 章 §2.2
\end{itemize}

% Auto-generated from ch01-basics.mdx — edit ch01-basics.mdx then re-run mdx2tex.mjs
\chapter{第 1 章 装饰器基础}
\label{ch:ch01-basics}

\begin{tcolorbox}[colback=blue!5,colframe=blue!40,title={📘 本章地图}]
本章解决三个问题：为什么装饰器能成立（一等函数 + 闭包）？\textbackslash\{\}textbackslash\textbackslash\{\}\{\textbackslash\{\}\}texttt\textbackslash\{\}\{@\textbackslash\{\}\} 到底是什么语法？如何写第一个装饰器？

\begin{itemize}
\item 1.1 函数是一等公民 - 1.2 闭包：内部函数记住环境 - 1.3 \textbackslash\{\}textbackslash\textbackslash\{\}\{\textbackslash\{\}\}textbackslash\textbackslash\{\}textbackslash\textbackslash\{\}\{\textbackslash\{\}\}\textbackslash\{\}\{\textbackslash\{\}textbackslash\textbackslash\{\}\{\textbackslash\{\}\}\textbackslash\{\}\}texttt\textbackslash\{\}textbackslash\textbackslash\{\}\{\textbackslash\{\}\}\textbackslash\{\}\{@\textbackslash\{\}textbackslash\textbackslash\{\}\{\textbackslash\{\}\}\textbackslash\{\}\} 语法糖：\textbackslash\{\}textbackslash\textbackslash\{\}\{\textbackslash\{\}\}textbackslash\textbackslash\{\}textbackslash\textbackslash\{\}\{\textbackslash\{\}\}\textbackslash\{\}\{\textbackslash\{\}textbackslash\textbackslash\{\}\{\textbackslash\{\}\}\textbackslash\{\}\}texttt\textbackslash\{\}textbackslash\textbackslash\{\}\{\textbackslash\{\}\}\textbackslash\{\}\{f = deco(f)\textbackslash\{\}textbackslash\textbackslash\{\}\{\textbackslash\{\}\}\textbackslash\{\}\} - 1.4 Recipe：计时器装饰器 - 1.5 常见陷阱 - 1.6 要点总结
\end{itemize}
\end{tcolorbox}
\section{1.1 函数是一等公民}
"一等公民"（first-class citizen）意味着函数可以：

\begin{enumerate}
\item 赋值给变量
\item 作为参数传给另一个函数
\item 作为另一个函数的返回值
\end{enumerate}

\begin{codeblock}{python}
# file: examples/first_class.py
def greet(name: str) -> str:
    """返回一条问候语。"""
    return f"Hello, {name}!"

# 1) 赋值给变量
say_hello = greet
print(say_hello("Alice"))  # Hello, Alice!

# 2) 作为参数传入高阶函数
def shout(fn, text: str) -> str:
    """调用 fn(text) 并把结果转为大写。"""
    return fn(text).upper()

print(shout(greet, "bob"))  # HELLO, BOB!
\end{codeblock}

这三条性质是所有装饰器的"地基"。没有一等函数，装饰器根本无从谈起。

\section{1.2 闭包：内部函数记住环境}
\textbackslash\{\}textbf\{闭包\}（closure）是指：一个内部函数引用了外部函数的局部变量，即使外部函数已经返回，内部函数仍然能访问这些变量。


\begin{codeblock}{python}
# file: examples/closure.py
def make_counter() -> callable:
    """返回一个闭包：每次调用计数 +1。"""
    count = 0  # 这个变量被闭包"记住"

    def increment() -> int:
        # nonlocal 告诉 Python：count 来自外层作用域
        nonlocal count
        count += 1
        return count

    return increment

# 调用：counter 就是 increment，却记住了 count
counter = make_counter()
print(counter())  # 1
print(counter())  # 2
print(counter())  # 3
\end{codeblock}

关键点：\textbackslash\{\}texttt\{counter()\} 每次都能读到自己那份 \textbackslash\{\}texttt\{count\}，即使 \textbackslash\{\}texttt\{make\_counter\} 早已执行完毕。装饰器就是"带额外行为的闭包"。

\section{1.3 \textbackslash\{\}texttt\{@\} 语法糖：\textbackslash\{\}texttt\{f = deco(f)\}}
PEP 318 引入的 \textbackslash\{\}texttt\{@\} 语法只是\textbackslash\{\}textbf\{语法糖\}（syntactic sugar）。下面两种写法\textbackslash\{\}textbf\{完全等价\}：


\begin{codeblock}{python}
# 写法 A：语法糖
@log
def hello(name):
    return f"hi {name}"

# 写法 B：手动调用
def hello(name):
    return f"hi {name}"
hello = log(hello)
\end{codeblock}

也就是说，装饰器是一个\textbackslash\{\}textbf\{输入为函数、输出为另一个函数\}的高阶函数：

\textbackslash\{\}[ T: f \textbackslash\{\}mapsto f' \textbackslash\{\}]

其中 \$f'\$ 通常是一个"包裹了 \$f\$ 的 wrapper"。

\subsection{最小可用装饰器}

\begin{codeblock}{python}
# file: examples/log_deco.py
import functools

def log(func):
    """打印函数调用日志的装饰器。"""

    # functools.wraps 拷贝原函数的 __name__、__doc__ 等元信息
    @functools.wraps(func)
    def wrapper(*args, **kwargs):
        print(f"[log] 调用 {func.__name__}, args={args}, kwargs={kwargs}")
        result = func(*args, **kwargs)  # 调用原函数
        print(f"[log] {func.__name__} 返回 {result!r}")
        return result

    return wrapper

@log
def add(a: int, b: int) -> int:
    return a + b

add(2, 3)
# [log] 调用 add, args=(2, 3), kwargs={}
# [log] add 返回 5
\end{codeblock}

图 1-1 是 \textbackslash\{\}texttt\{@log\} 装饰器的调用时序：

\textbackslash\{\}begin\{figure\}[htbp]\textbackslash\{\}centering\textbackslash\{\}includegraphics[width=0.85\textbackslash\{\}textwidth]\{figures/decorator-sequence.svg\}\textbackslash\{\}caption\{图 1-1：@log 装饰器调用时序\}\textbackslash\{\}end\{figure\}

如图所示，每次调用 \textbackslash\{\}texttt\{add(2, 3)\} 实际上进入的是 \textbackslash\{\}texttt\{wrapper\}，它先打日志、再调原函数、最后返回结果。

\subsection{装饰器 vs 普通函数}
\begin{table}[htbp]
\centering
\begin{tabularx}{\textwidth}{X X X}
\toprule
维度 & 普通函数 & 装饰器 \\
\midrule
输入 & 数据 & \textbackslash\{\}textbf\{另一个函数\} \\
输出 & 数据 & \textbackslash\{\}textbf\{包裹后的函数\} \\
用法 & \textbackslash\{\}texttt\{result = f(x)\} & \textbackslash\{\}texttt\{f = deco(f)\}（或 \textbackslash\{\}texttt\{@deco\}） \\
典型场景 & 业务逻辑 & 横切关注点：日志/缓存/权限 \\
\bottomrule
\end{tabularx}
\end{table}

\section{1.4 Recipe：计时器装饰器}
<div className="callout recipe">

\subsubsection{🍳 Recipe 1-1：测量函数耗时}
\textbackslash\{\}textbf\{Problem\}：想知道一个函数运行了多久，又不想把计时代码散落到每处调用。

\textbackslash\{\}textbf\{Solution\}：


\begin{codeblock}{python}
# file: recipes/timer.py
import time
import functools

def timer(func):
    @functools.wraps(func)
    def wrapper(*args, **kwargs):
        t0 = time.perf_counter()       # 开始计时（高精度）
        try:
            return func(*args, **kwargs)
        finally:
            dt = time.perf_counter() - t0
            print(f"[timer] {func.__name__} 耗时 {dt*1000:.2f} ms")
    return wrapper

@timer
def slow_sum(n: int) -> int:
    return sum(range(n))

slow_sum(10_000_000)  # [timer] slow_sum 耗时 253.47 ms
\end{codeblock}

\textbackslash\{\}textbf\{Discussion\}：

\begin{itemize}
\item 用 \textbackslash\{\}texttt\{try/finally\} 保证即使函数抛异常也会打印耗时。
\item \textbackslash\{\}texttt\{time.perf\_counter()\} 是 Python 推荐的短时计时 API（比 \textbackslash\{\}texttt\{time.time()\} 精度更高，且不受系统时钟回拨影响）。
\end{itemize}
\textbackslash\{\}textbf\{See Also\}：第 2 章将把它升级为带参数的 \textbackslash\{\}texttt\{@timer(unit="ms")\}。

</div>

\section{1.5 常见陷阱}
<div className="callout pitfall">

\subsubsection{⚠️ 陷阱 1：忘记 \textbackslash\{\}texttt\{@functools.wraps(func)\}}
如果不加 \textbackslash\{\}texttt\{@functools.wraps(func)\}，装饰后的函数会丢失原函数的元信息：


\begin{codeblock}{python}
@log
def add(a, b):
    """两个整数相加。"""
    return a + b

print(add.__name__)  # 没有 wraps → "wrapper"，加了 wraps → "add"
print(add.__doc__)   # 没有 wraps → None，加了 wraps → "两个整数相加。"
\end{codeblock}

这会破坏调试器、文档生成工具（Sphinx）、以及依赖 \textbackslash\{\}texttt\{\_\_name\_\_\} 的框架（如 Flask 路由注册）。

\textbackslash\{\}textbf\{正确做法\}：\textbackslash\{\}textbf\{永远\}在 wrapper 上加 \textbackslash\{\}texttt\{@functools.wraps(func)\}。

</div>

\section{1.6 要点总结}
\begin{tcolorbox}[colback=green!5,colframe=green!40,title={✅ 核心要点}]
\begin{itemize}
\item 装饰器的前提是\textbackslash\{\}textbackslash\textbackslash\{\}\{\textbackslash\{\}\}textbf\textbackslash\{\}\{函数是一等公民\textbackslash\{\}\} + \textbackslash\{\}textbackslash\textbackslash\{\}\{\textbackslash\{\}\}textbf\textbackslash\{\}\{闭包\textbackslash\{\}\}。 - \textbackslash\{\}textbackslash\textbackslash\{\}\{\textbackslash\{\}\}texttt\textbackslash\{\}\{@deco\textbackslash\{\}\} 只是 \textbackslash\{\}textbackslash\textbackslash\{\}\{\textbackslash\{\}\}texttt\textbackslash\{\}\{f = deco(f)\textbackslash\{\}\} 的语法糖。 - 一个合格的装饰器必须用 \textbackslash\{\}textbackslash\textbackslash\{\}\{\textbackslash\{\}\}texttt\textbackslash\{\}\{@functools.wraps(func)\textbackslash\{\}\} 保留原函数元信息。 - wrapper 用 \textbackslash\{\}textbackslash\textbackslash\{\}\{\textbackslash\{\}\}texttt\textbackslash\{\}\{\textbackslash\{\}textbackslash\textbackslash\{\}\{\textbackslash\{\}\}emph\textbackslash\{\}\{args, \textbackslash\{\}textbackslash\textbackslash\{\}\{\textbackslash\{\}\}textbf\textbackslash\{\}\{kwargs\textbackslash\{\}\} 才能适配任意签名的被装饰函数。 - 装饰器是\textbackslash\{\}\}横切关注点\textbackslash\{\}\}*（日志、计时、权限、缓存）的首选表达。
\end{itemize}
\end{tcolorbox}
\section{1.7 延伸阅读}
\begin{itemize}
\item 第 2 章 · 进阶装饰器：带参装饰器、类装饰器、堆叠顺序
\item 附录 A · functools 速查
\item PEP 318 – Decorators for Functions and Methods：https://peps.python.org/pep-0318/
\end{itemize}
---

\textbackslash\{\}textbf\{本章参考文献\}

\begin{enumerate}
\item Python Software Foundation. "PEP 318 – Decorators for Functions and Methods." https://peps.python.org/pep-0318/
\item Real Python. "Primer on Python Decorators." https://realpython.com/primer-on-python-decorators/
\end{enumerate}

% Auto-generated from ch02-advanced.mdx — edit ch02-advanced.mdx then re-run mdx2tex.mjs
\chapter{第 2 章 参数化装饰器与类装饰器}
\label{ch:ch02-advanced}

\begin{tcolorbox}[colback=blue!5,colframe=blue!40,title={📘 本章地图}]
第 1 章的装饰器只能"裸用"（\textbackslash\{\}textbackslash\textbackslash\{\}\{\textbackslash\{\}\}texttt\textbackslash\{\}\{@log\textbackslash\{\}\}）。本章学会带参数的装饰器、装饰类的装饰器，以及多个装饰器堆叠时的执行顺序。

\begin{itemize}
\item 2.1 带参装饰器：为什么多一层函数 - 2.2 Recipe：@retry(times=N) - 2.3 类装饰器：装饰整个类 - 2.4 装饰器堆叠：自下而上应用 - 2.5 常见陷阱 - 2.6 要点总结
\end{itemize}
\end{tcolorbox}
\section{2.1 带参装饰器：为什么多一层函数}
第 1 章的 \textbackslash\{\}texttt\{@log\} 写死了行为。如果想 \textbackslash\{\}texttt\{@retry(times=3)\} 这样\textbackslash\{\}textbf\{带参数\}，必须多包一层函数：


\begin{codeblock}{python}
# 不带参：@deco           → deco(func) 返回 wrapper
# 带参：  @deco(times=3)  → deco(times=3)(func) 返回 wrapper
\end{codeblock}

也就是说，带参装饰器是\textbackslash\{\}textbf\{返回真正装饰器的工厂\}。通用模板：


\begin{codeblock}{python}
def deco_factory(params):
    """外层：接收装饰器参数，返回真正的装饰器。"""
    def real_decorator(func):
        """中层：真正的装饰器，输入 func，输出 wrapper。"""
        @functools.wraps(func)
        def wrapper(*args, **kwargs):
            """内层：被装饰后的函数。"""
            ...  # 可以使用 params、func
        return wrapper
    return real_decorator
\end{codeblock}

函数嵌套三层，理解关键是：\textbackslash\{\}textbf\{每次 \textbackslash\{\}texttt\{@deco(...)\} 先执行 \textbackslash\{\}texttt\{deco(...)\}，把结果当装饰器用。\}

\section{2.2 Recipe：@retry(times=N)}
<div className="callout recipe">

\subsubsection{🍳 Recipe 2-1：失败自动重试}
\textbackslash\{\}textbf\{Problem\}：调用不稳定的外部 API 时，想自动重试 N 次再放弃。

\textbackslash\{\}textbf\{Solution\}：


\begin{codeblock}{python}
# file: recipes/retry.py
import time
import functools

def retry(times: int = 3, delay: float = 0.5):
    """失败重试装饰器。

    Args:
        times: 最大重试次数（含首次调用）。
        delay: 每次重试前等待的秒数，避免雪崩。
    """
    def decorator(func):
        @functools.wraps(func)
        def wrapper(*args, **kwargs):
            last_exc = None
            for attempt in range(1, times + 1):
                try:
                    return func(*args, **kwargs)
                except Exception as exc:
                    last_exc = exc
                    print(f"[retry] 第 {attempt}/{times} 次失败：{exc}")
                    if attempt < times:
                        time.sleep(delay)  # 退避等待
            # 全部失败，向外抛出最后一次异常
            raise last_exc
        return wrapper
    return decorator

import random

@retry(times=3, delay=0.1)
def flaky_api() -> str:
    if random.random() < 0.7:
        raise ConnectionError("网络抖动")
    return "ok"

print(flaky_api())
\end{codeblock}

\textbackslash\{\}textbf\{Discussion\}：

\begin{itemize}
\item 三层嵌套对应 \textbackslash\{\}texttt\{retry(times=3)(func)(args)\}。
\item 用 \textbackslash\{\}texttt\{time.sleep(delay)\} 做简单退避；生产环境建议\textbackslash\{\}textbf\{指数退避\}（exponential backoff）。
\item 最后一次重试失败才 \textbackslash\{\}texttt\{raise\}，中间失败只打印日志。
\end{itemize}
</div>

\section{2.3 类装饰器：装饰整个类}
装饰器不仅能装饰函数，也能装饰\textbackslash\{\}textbf\{类\}。最常见的用途：

\begin{itemize}
\item 给类自动注册到某个全局表（如 ORM 模型、插件系统）
\item 把类改造为单例（singleton）
\item 自动给类加方法/属性
\end{itemize}

\begin{codeblock}{python}
# file: examples/singleton.py
import functools

def singleton(cls):
    """把类改造成单例：cls() 永远返回同一个实例。"""
    _instances = {}

    @functools.wraps(cls, updated=())
    def wrapper(*args, **kwargs):
        if cls not in _instances:
            _instances[cls] = cls(*args, **kwargs)
        return _instances[cls]

    return wrapper

@singleton
class AppConfig:
    def __init__(self):
        self.env = "production"

a = AppConfig()
b = AppConfig()
assert a is b  # True，同一个实例
\end{codeblock}

图 2-1 展示了装饰器相关类之间的继承与组合关系：

\textbackslash\{\}begin\{figure\}[htbp]\textbackslash\{\}centering\textbackslash\{\}includegraphics[width=0.85\textbackslash\{\}textwidth]\{figures/decorator-class.svg\}\textbackslash\{\}caption\{图 2-1：装饰器类层次\}\textbackslash\{\}end\{figure\}

注意：类装饰器接收的参数是 \textbackslash\{\}textbf\{类对象本身\}（而不是类的某个实例）。

\section{2.4 装饰器堆叠：自下而上应用}
多个装饰器可以堆叠在同一个函数上：


\begin{codeblock}{python}
@timer
@retry(times=3)
@log
def fetch(url): ...
\end{codeblock}

\textbackslash\{\}textbf\{执行顺序\}规则只有一句话：\textbackslash\{\}textbf\{离函数最近的装饰器最先应用（自下而上）\}。

上面的例子等价于：


\begin{codeblock}{python}
fetch = timer(retry(times=3)(log(fetch)))
\end{codeblock}

应用顺序是 \textbackslash\{\}texttt\{log → retry → timer\}；而\textbackslash\{\}textbf\{执行顺序\}（调用 \textbackslash\{\}texttt\{fetch()\} 时 wrapper 的堆叠顺序）正好相反：\textbackslash\{\}texttt\{timer → retry → log → 原 fetch\}。

记忆法：\textbackslash\{\}textbf\{装饰像洋葱，外到内层层包裹\}。

\subsection{装饰器类型速查}
\begin{table}[htbp]
\centering
\begin{tabularx}{\textwidth}{X X X X}
\toprule
类型 & 语法 & 层数 & 典型场景 \\
\midrule
无参函数装饰器 & \textbackslash\{\}texttt\{@deco\} & 2 层 & \textbackslash\{\}texttt\{@log\}, \textbackslash\{\}texttt\{@timer\} \\
带参函数装饰器 & \textbackslash\{\}texttt\{@deco(x=1)\} & 3 层 & \textbackslash\{\}texttt\{@retry(times=3)\}, \textbackslash\{\}texttt\{@lru\_cache(maxsize=64)\} \\
类装饰器 & \textbackslash\{\}texttt\{@deco\} 在 \textbackslash\{\}texttt\{class\} 上 & 2 层 & \textbackslash\{\}texttt\{@singleton\}, \textbackslash\{\}texttt\{@dataclass\} \\
堆叠装饰器 & 多个 \textbackslash\{\}texttt\{@\} & 多套 & 组合日志+重试+计时 \\
\bottomrule
\end{tabularx}
\end{table}

\section{2.5 常见陷阱}
<div className="callout pitfall">

\subsubsection{⚠️ 陷阱 2：装饰器堆叠顺序写反导致行为错乱}

\begin{codeblock}{python}
@log        # 外层
@retry(3)   # 内层
def api(): ...
\end{codeblock}

如果你期望"每次重试都打日志"，这个顺序是对的（外层 log 包裹整个 retry 循环，只会打一次）。如果你期望"每次尝试都打日志"，必须把 \textbackslash\{\}texttt\{@log\} 写在 \textbackslash\{\}texttt\{@retry\} \textbackslash\{\}textbf\{里面\}：


\begin{codeblock}{python}
@retry(3)
@log
def api(): ...
\end{codeblock}

\textbackslash\{\}textbf\{规则\}：想让某个 wrapper 对"每次重试"生效，就让它更靠近被装饰函数。

\subsubsection{⚠️ 陷阱 3：类装饰器忘记保留类元信息}
装饰类时 \textbackslash\{\}texttt\{functools.wraps(cls)\} 默认不会复制类的 \textbackslash\{\}texttt\{\_\_doc\_\_\}/\textbackslash\{\}texttt\{\_\_name\_\_\}，需要加 \textbackslash\{\}texttt\{updated=()\} 参数：


\begin{codeblock}{python}
@functools.wraps(cls, updated=())
def wrapper(*args, **kwargs): ...
\end{codeblock}

</div>

\section{2.6 要点总结}
\begin{tcolorbox}[colback=green!5,colframe=green!40,title={✅ 核心要点}]
\begin{itemize}
\item 带参装饰器 = "返回装饰器的工厂"，\textbackslash\{\}textbackslash\textbackslash\{\}\{\textbackslash\{\}\}textbf\textbackslash\{\}\{三层嵌套\textbackslash\{\}\}：工厂 → 装饰器 → wrapper。 - 类装饰器接收\textbackslash\{\}textbackslash\textbackslash\{\}\{\textbackslash\{\}\}textbf\textbackslash\{\}\{类对象\textbackslash\{\}\}，常见用途是单例、注册、改造类。 - 多个装饰器堆叠：\textbackslash\{\}textbackslash\textbackslash\{\}\{\textbackslash\{\}\}textbf\textbackslash\{\}\{自下而上\textbackslash\{\}\}应用（离函数最近最先执行装饰），运行时\textbackslash\{\}textbackslash\textbackslash\{\}\{\textbackslash\{\}\}textbf\textbackslash\{\}\{外→内\textbackslash\{\}\}包裹。 - 用 \textbackslash\{\}textbackslash\textbackslash\{\}\{\textbackslash\{\}\}texttt\textbackslash\{\}\{functools.wraps\textbackslash\{\}\} 保留元信息；装饰类时加 \textbackslash\{\}textbackslash\textbackslash\{\}\{\textbackslash\{\}\}texttt\textbackslash\{\}\{updated=()\textbackslash\{\}\}。
\end{itemize}
\end{tcolorbox}
行内公式示例：装饰器的数学抽象是 \$T: f \textbackslash\{\}mapsto f'\$，其中 \$f'\$ 是 wrapper。

块级公式示例：带参装饰器的三次调用可以写成：

\textbackslash\{\}[ D(a, b)(f)(x) = W\_\{a,b\}\textbackslash\{\}bigl(f(x)\textbackslash\{\}bigr) \textbackslash\{\}]

\section{2.7 延伸阅读}
\begin{itemize}
\item 第 1 章 · 装饰器基础
\item 附录 A · functools 速查
\item Python 官方文档 \textbackslash\{\}texttt\{functools\}：https://docs.python.org/3/library/functools.html
\end{itemize}
---

\textbackslash\{\}textbf\{本章参考文献\}

\begin{enumerate}
\item Python Software Foundation. "functools — Higher-order functions and operations on callable objects." https://docs.python.org/3/library/functools.html
\end{enumerate}

% Auto-generated from index.mdx — edit index.mdx then re-run mdx2tex.mjs
\chapter{前言}
\label{ch:index}

\section*{Python 装饰器 CookBook}
\begin{tcolorbox}[colback=blue!5,colframe=blue!40,title={📘 本册地图}]
本 CookBook 面向中级 Python 开发者，系统讲解装饰器（Decorator）这一元编程核心工具。

\begin{itemize}
\item 第 1 章 · 装饰器基础：闭包、\textbackslash\{\}textbackslash\textbackslash\{\}\{\textbackslash\{\}\}texttt\textbackslash\{\}\{@\textbackslash\{\}\} 语法糖、无参装饰器 - 第 2 章 · 进阶装饰器：带参装饰器、类装饰器、装饰器堆叠 - 附录 A · functools 速查
\end{itemize}
\end{tcolorbox}
\section{为什么要学装饰器}
装饰器是 Python 中最具"Pythonic"风格的语言特性之一。从 Web 框架（Flask 的 \textbackslash\{\}texttt\{@app.route\}、FastAPI 的 \textbackslash\{\}texttt\{@app.get\}）到权限校验、日志记录、缓存、重试，几乎所有横切关注点（cross-cutting concern）都可以用装饰器优雅地表达。

理解装饰器不仅是会写 \textbackslash\{\}texttt\{@xxx\}，更要理解：

\begin{itemize}
\item \textbackslash\{\}textbf\{函数是一等公民\}（first-class citizen）：函数可以作为参数、返回值、赋值给变量。
\item \textbackslash\{\}textbf\{闭包\}（closure）：内部函数记住外部函数的词法环境。
\item \textbackslash\{\}textbf\{语法糖只是函数调用\}：\textbackslash\{\}texttt\{@deco\} 本质上是 \textbackslash\{\}texttt\{f = deco(f)\}。
\end{itemize}
\section{全景图}

\begin{tcolorbox}[colback=gray!5,colframe=gray!40,title={Diagram: mermaid}]
Diagram source kept in \texttt{diagrams/}. Rendered in Nextra/PDF.\end{tcolorbox}

> 上面这张心智图（mindmap）覆盖了本册的全部知识点。建议先通览，再进入第 1 章。

\section{目标读者}
\begin{itemize}
\item 已经会写 Python 函数与类，读过至少一份 Python 项目源码
\item 想系统理解 \textbackslash\{\}texttt\{@decorator\} 的内部机理，而不是只会复制粘贴
\item 目标：能写出可复用、可维护、可调试的装饰器
\end{itemize}
\section{约定}
\begin{itemize}
\item 所有示例基于 \textbackslash\{\}textbf\{Python 3.10+\}（使用 PEP 604 的 \textbackslash\{\}texttt\{X | Y\} 类型注解语法）。
\item 代码块标注文件路径，可直接复制运行。
\item 每章末尾的 \textbackslash\{\}textbf\{Recipes\} 提供任务导向示例；\textbackslash\{\}textbf\{Pitfalls\} 列出常见坑。
\end{itemize}


\end{document}
